\documentclass[12pt]{article}
\usepackage[utf8]{inputenc}
\usepackage[T1]{fontenc}
\usepackage{hyperref}
\title{Stoney Language Resources for K-12 Curriculum Development}
\author{Compiled by Claude Sonnet 5.5 (claude-sonnet-5-5) for Andr\'as Kornai}
\date{}
\begin{document}
\maketitle

\section{Introduction}
Stoney (Stoney Nakoda, ISO 639-3 \texttt{sto}) is a Siouan language of Alberta, Canada. It is spoken by the Bearspaw, Chiniki and Wesley First Nations, about 70 km west of Calgary. About 3,200 speakers are reported. Although the entry says it is not known to be taught in schools, the searches found school materials produced by the Stoney Education Authority.

\section{Primary Educational Resources}
\begin{itemize}
\item In January 2023 the Stoney Education Authority (SEA) released a Stoney Nakoda Level 2 textbook, a printed Stoney Nakoda Student Dictionary and a podcast of Elders telling stories (CTV News): \url{https://calgary.ctvnews.ca/new-textbook-podcast-aim-to-keep-stoney-nakoda-language-alive-in-alberta-1.6242183}
\item The first set of materials (2021) was aimed at elementary students. It included a Level 1 textbook and picture books (Global News): \url{https://globalnews.ca/news/9430501/stoney-nakota-language-textbook/}
\item CBC coverage of the Stoney Nakoda language effort: \url{https://amp.cbc.ca/news/canada/calgary/stoney-nakoda-nation-language-alberta-1.6723532}
\end{itemize}

\section{Linguistic and Reference Materials}
\begin{itemize}
\item The Stoney Nakoda Student Dictionary has about 9,000 words. It was first available as an online and mobile dictionary and was later printed. The Language Conservancy report describes how it was made: \url{https://languageconservancy.org/new-textbook-podcast-aim-to-keep-stoney-nakoda-language-alive-in-alberta/}
\end{itemize}

\section{Religious and Community Resources}
The materials were made with Elders from the three First Nations (Bearspaw, Chiniki and Wesley), with support workers and linguists. The searches found no scripture resources.

\section{Digital and Audio Resources}
\begin{itemize}
\item The online and mobile dictionary and the Elders' story podcast, both described in the coverage above. The searches did not return their direct URLs. The articles point to stoneyeducation.ca, which was not itself seen in search results.
\end{itemize}

\section{Curriculum Development Context}
The Language Conservancy and the SEA used a rapid word collection method to build the dictionary in about 10 days of intensive sessions. Stoney is still a first language for many adults on the three reserves, and about 1,500 students on the First Nation use a basic Stoney textbook. Level 1 and Level 2 textbooks exist, so higher levels and materials for other subjects are the obvious gaps. Current availability and pricing were not found, so the SEA should be contacted directly.

\end{document}